\section{C Programming}
\subsubsection{von Neumann Architecture}
\begin{figure}[h!]
  \centering
  \includegraphics[width=0.4\textwidth]{vonNeumannArchi}
\end{figure}
The von Neumann architecture (also known as the \textit{Princeton
architecture}, or \textit{stored-program architecture}) describes a computer
consisting of:
\begin{itemize}
  \item Central Processing Unit (CPU) - Registers, A control unit containing an instruction register and program counter, An arithmetic/logic unit (ALU)
  \item Memory - Stores both program and data in random-access memory (RAM)
  \item I/O Devices
\end{itemize}

\subsubsection{Program Structure}
The form of a C Program typically looks like the following:
\begin{lstlisting}
preprocessor directives

structure types

function prototypes

main function header {
  declaration of variables
  executable statements
}

function definitions
\end{lstlisting}
\begin{figure}[h!]
  \centering
  \includegraphics[width=0.8\textwidth]{cProgram}
\end{figure}

\begin{remark}
  The main function is obligatory in C. Typically, ``executable statements''
  consists 3 parts: \textit{Input data, Computation, Output Results}.
\end{remark}
\paragraph{Function Prototypes}
User-defined function prototypes are of the syntax
\begin{lstlisting}
  return_type function_name(parameter_type);
\end{lstlisting}
The names of parameters in the function prototype are optional and ignored by
the compiler.
\begin{remark}
 It is good practice to put function prototypes before the main function. 
\end{remark}
If no prototype is provided but the function is used, the compiler adds an
\textbf{implicit declaration} and returns a warning. The default return type is
assumed to be \verb|int|, and hence is almost always wrong.
\begin{example}
  The function parameter types may consist of arrays, strings and structs as well.
  \begin{lstlisting}
    int func(int *, int [], int);
  \end{lstlisting}
\end{example}

\subsection{Preprocessor Directives}
The C preprocessor provides the following:
\begin{itemize}
  \item Inclusion of header files:
    \begin{itemize}
      \item To use input/output functions such as \verb|scanf()| and
        \verb|printf()|: \verb|#include <stdio.h>| (compile with the \verb|-o|
        flag to specify output file)
      \item To use mathematical functions: \verb|#include <math.h>| and compile
        with the \verb|-lm| flag on sunfire.
      \item To use string functions: \verb|#include <string.h>|
    \end{itemize}
  \item Macro expansions:
    To define a macro for a constant value. Useful for declaring constants, to
    use all CAPS. \verb|#define PI 3.142| (without a semicolon, otherwise it
    will be included in the substitution).
  \item Conditional Compilation (Will not be introduced in CS2100)
\end{itemize}
\subsection{Input / Output}
We use format specifiers that serve as placeholders for values to be displayed
or read.
\vspace{2mm}

\centering
\begin{tabular}{c c c}
  \toprule
  Placeholder & Variable Type & Function Use \\
  \midrule
  \verb|%c| & char & printf / scanf \\
  \verb|%d| & int & printf / scanf \\
  \verb|%f| & float / double & printf \\
  \verb|%f| & float & scanf \\
  \verb|%lf|  & double & scanf \\
  \verb|%e| & float / double & printf (with scientific notation) \\
  \verb|%p| & address & printf (in hexadecimal)\\
  \bottomrule
\end{tabular}

\raggedright
\begin{example}[Format Specifiers]
  Examples of format specifiers used in \verb|printf()| are:
  \begin{itemize}
    \item \verb|%5d|: displays an integer in a width of 5, right justified.
    \item \verb|%8.3f|: displays a real number (float or double) in a width of 8,
      with 3 decimal places, right justified.
  \end{itemize}
  For \verb|scanf()|, format specifiers are used without indicating width or
  decimal places.
\end{example}

Additionally, we may use escape sequences in \verb|printf()| for certain special
effects or to display certain characters properly.
\vspace{2mm}

\centering
\begin{tabular}{c c c}
  \toprule
  Escape Sequence & Meaning & Result \\
  \midrule
  \verb|\n| & New line & Subsequent output appears on the next line \\
  \verb|\t| & Horizontal tab & Move to the next tab position on the current line \\
  \verb|\"| & Double quote & Displays a double quote " \\
  \verb|%%| & Percent & Dispalys a percent character \% \\
  \bottomrule
\end{tabular}

\raggedright
\subsection{Computation}
\subsubsection{Functions}
Computation in C is done through functions. A function body mainly composes of: 
\begin{itemize}
  \item \textit{Declaration statements} introducing an identifier and
    describing its type. It is required by the compiler to accept references to
    that identifier. Re-declaration is legal outside of class scope.
    \begin{remark}
      \verb|extern| can be used to import variables from another file. On
      variables, the variable is not allocated any memory. On functions, extern
      is the default and otherwise require \verb|static|. Similarly, using
      \verb|static| on a global variable ensures that it cannot be named
      outside of the current file. 
    \end{remark}
  \item \textit{Definition Statements} instantiating or implementing the identifiers. It
    is required by the linker to link references. It can be used in place of
    declarations.
  \item \textit{Executable statements} describing the processing on the memory
    cells.
\end{itemize}
\begin{remark}
  Using \verb|static| \textit{on a local variable} persists the value of the
  variable between function calls.
\end{remark}
\begin{remark}
  The special \verb|int main(void)| function is the entry point of the program.
\end{remark}

In C, function parameters are passed to the formal parameters through
\textit{pass-by-value}. That is, values of the actual parameters are copied to
the formal parameters.

\paragraph{Scope Rule}
\begin{definition}[Lexical Scoping]
  C uses \textit{lexical scoping}, also known as static scoping. That is, the
  compiler is able to deduce the scope of each variable at compile-time.
\end{definition}
Formal parameters and variables declared in the function are local to the
function. These local variables are only accessible in the function they are
declared. This is known as \textit{scope rule}. It is implemented as follows:
\begin{itemize}
  \item When a function is called, an activation record is created on the call
    stack, and memory is allocated for local parameters and variables (known as
    \textit{automatic variables}).
  \item When the function terminates, the activation record is removed, and
    memory allocated for the local parameters and variables is released.
\end{itemize}
In contrast to automatic variables, we have static variables that remain in
memory after the function is executed.

\subsubsection{Variables}
Data used in a program are stored in \textbf{variables}, identified by a name
(identifier), having a data type, containing a value which can be modified and
having an address (which varies from run to run).
\begin{remark}
  A variable is declared with a data type, and may be initialised during declaration
  \begin{lstlisting}
    int count = 3;
  \end{lstlisting}
  Without initialisation, the variable contains an unknown value (\textit{not
  necessarily zero}). We can use the \verb|-Wall| option during compilation to
  turn on warnings for missing initialisation.
\end{remark}
\begin{lstlisting}
int x;
scanf("%d", &x); // Refers to the (address of) the memory cell of x
printf("x equals %d\n", x) // Refers to the value in the variable x
\end{lstlisting}
C is a relatively strongly typed language, where every variable is to be
declared with a data type, but less so than Java or Pascal, as C supports
implicit type conversions and allows pointer values to be explicitly cast.
\subsubsection{Data Types}
In C, we have the basic data types (sizes in \textit{sunfire}):
\begin{itemize}
  \item int: $4\text{ bytes}; -2,147,483,648\ (-2^{31})\text{ to }+2,147,483,647\ (2^{31} – 1)$
  \item float: $4\text{ bytes};$ or double: $8\text{ bytes}$; May use
  scientific notation (e.g. \verb|1.5e-2|)
  \item char: Enclosed in a pair of single quotes.
\end{itemize}
The size of a type, or a variable, can be found using the \verb|sizeof()| function.\\
\begin{remark}
  With \verb|stdint.h|, we have more precise/unsigned integer types
  \verb|int8_t|, \verb|uint8_t|, \verb|int16_t|, \verb|uint16_t|,
  \verb|int32_t|, \verb|uint32_t|, \verb|int64_t|, \verb|uint64_t|, which can
  be printed with \verb|inttypes.h|.
\end{remark}
\subsubsection{User-defined Identifiers}
To define functions, or declare the use of variables, the user is required to
define the identifier as the name of a variable or function. It is obligatory
to follow the following rules:
\begin{itemize}
  \item Consists of letters (a-z, A-Z), digits (0-9), and underscores (\_)
    but cannot begin with a digit.
  \item Cannot be a reserved word. See \url{http://c.ihypress.ca/reserved.html}
    for a complete list.
\end{itemize}
In addition, it is recommended (while not obligatory) to follow these restrictions:
\begin{itemize}
  \item Begins with a lowercase letter, except constants which we identify in ALL CAPS.
  \item Avoids standard identifiers (names of common functions) such as
  \verb|printf|, \verb|scanf|.
  \begin{lstlisting}
    # include <stdio.h>
    int printf=5 // No issue, if printf() is not used
    printf("Hello world!") // Compile time error thrown from this line.
  \end{lstlisting}
\end{itemize}
\subsubsection{Assignment Statements}
Executable statements, in general, are I/O statements or computational and
assignment statements. Assignment statements store a value in a variable:
\verb|lvalue = 30;|.
\begin{remark}
  \verb|lvalue| must be assignable. That is, \verb|32 = a| (not a variable)
  and \verb|a + b = c| (an expression) are not assignable and throws the
  compilation error "lvalue required as left operand of assignment".
\end{remark}
Additionally, assignment can be cascaded, with associativity from right to left:

\centering
\verb|a = b = c = 1 + 2| $\equiv$ \verb|a = (b = (c = 1 + 2))|

\raggedright
In C, the assignment statement has the side effect of returning the value of
the expression. This allows the cascading of assignment as shown, but also
allows statements such as \verb|a = 5 + (b = 10)|.
\subsubsection{Operations}
\begin{tabular}{c c c}
  \toprule
  Operator & Operator & Associativity \\
  \midrule
  Primary Expression Operators & \verb|()  []  .  ->  expr++  expr--| & Left to right \\
  Unary Operators & \verb|* & + ! ~ - ++expr --expr (typecast) sizeof| & Right to left \\
  \multirow{6}{8em}{Binary Operators} & \verb|*  /  %| & \multirow{6}{5.3em}{Left to right} \\
   & \verb|+  -| &  \\
   & \verb|<  >  <=  >=| &  \\
   & \verb|==  !=| &  \\
   & \verb|&&| &  \\
   & \verb$||$ &  \\
  Ternary Operator & \verb|predicate ? consequent : alternative| & Right to left\\
  Assignment Operators & \verb|=  +=  -=  *=  /=  %=| & Right to left\\
  \bottomrule
\end{tabular}
\begin{remark}
  In C, the operator \verb|%| calculates the remainder \verb|a%b| as follows:
  \[\left(a/b\right)*b+a\ \%\ b=a.\] 
\end{remark}
\begin{remark}
  There is no Boolean type in ANSI C. We use \verb|0| to represent
  \textit{false}, and any other value to represent \text{true}.
\end{remark}
Mixed-typed arithmetic operations:
\begin{multicols}{2}
\begin{itemize}
  \item \verb|int  m = 10/4; // m = 2|
  \item \verb|float p = 10/4; // p = 2.0|
  \item \verb|int n = 10/4.0; // n = 2|
  \item \verb|float q = 10/4.0; // q = 2.5|
\end{itemize}
\begin{itemize}
  \item \verb|int r = -10/4.0; // r = -2|
  \item \verb|float pp = (float) aa/4; // pp = 1.5|
  \item \verb|int nn = (int) ff / aa; // nn = 2|
  \item \verb|int qq = (float) (aa/4); // qq = 1.0|
\end{itemize}
\end{multicols}

\subsubsection{Math Library}
With inclusion of the \verb|<math.h>| library, compiled with the \verb|-lm|
option, we are able to use the following common math functions:
\begin{multicols}{2}
  \begin{itemize}
    \item \verb|int abs(int x)|
    \item \verb|double fabs(double x)|
    \item \verb|double ceil(double x)|
    \item \verb|double floor(double x)|
    \item \verb|double exp(double x)|
    \item \verb|double pow(double x, double y)|
  \end{itemize}
  \begin{itemize}
    \item \verb|double log(double x)|
    \item \verb|double log10(double x)|
    \item \verb|double sqrt(double x)|
    \item \verb|double cos(double x)| (radians)
    \item \verb|double sin(double x)| (radians)
    \item \verb|double tan(double x)| (radians)
  \end{itemize}
\end{multicols}

\subsection{Pointers}
Using pointers, we are able to directly manipulate memory contents. We are able
to modify the value of variables outside of a function's scope within the
function, by passing in the pointer into the function.
\begin{definition}[Pointer]
  A variable that contains the address of another variable is known as a
  pointer, or a pointer variable.
\end{definition}
To work with addresses, we have the unary \textit{address operator} \verb|&| to
obtain the address of a given variable and \textit{dereferencing/indirection
operator} \verb|*| to access the value of a given address.

To declare a pointer to a variable of \verb|type|, we use the syntax
\begin{lstlisting}
  type *variable_ptr; // Convention to name pointer with suffix _ptr or _p
\end{lstlisting}

\begin{remark}
  Using the address of an uninitialised variable is well-defined behaviour, but using
  a uninitialised pointer is undefined behaviour.
  \begin{lstlisting}
    char x;
    printf("%p\n", &x); // Prints address of x
    char *y;
    printf("%p\n", y); // Issues warning
  \end{lstlisting}
\end{remark}

Using pointers can lead to issues relating to correctness and security, such as:
\begin{itemize}
  \item \textbf{Aliasing}: When two pointers point to the same object and we
    might not be able to tell.
  \item \textbf{Invalid/NULL pointer dereferencing}: Inability to check validity of a pointer.
  \item \textbf{Dangling Pointers}: When the memory object being pointed has
    been freed/is no longer available.
\end{itemize}

\begin{remark}
  Since all pointers have the same size and can point to the storage of no
  type, we can create a \textit{generic pointer}
  \[\verb|void *ptr|.\]
   This allows for a pointer that refers to a value that has unknown type, or
   several possible types, and is often used as a parameter of generic
   functions. They should be \textit{typecast} before valid usage.
\end{remark}

\subsubsection{Pointer Arithmetic}
When a pointer is incremented, the address is incremented by the size of the
pointer's type. For instance, a int pointer will be incremented by 4 bytes,
while a char pointer will be incremented by 1 byte.

\subsubsection{Dynamic Storage Allocation}
During runtime, we may allocate storage by using the \verb|void *malloc(size)|
system call. It exists as a pair to the \verb|free((void *)ptr)| system call to
free the pointer. However, this requires caution as it may result in dangling
pointers and hence undefined behaviour.

\subsection{Arrays, Strings and Structures}
To further organise data beyond the basic data types, with more logical
representations and greater ease of manipulations, we may use arrays, strings
and structs in C.
\subsubsection{Arrays}
Arrays are defined by the \textit{element type, array name and maximum size}. They
occupy \textbf{contiguous} memory locations and are accessed through indexing.
The array name itself is a pointer, which refers to the address of the
first element of that array. That is,
\[\verb|arr|\equiv \verb|&arr[0]|.\] 
Additionally, it is a \textbf{fixed pointer}. It is hence \textit{illegal} to
alter the value of this variable.
\begin{remark}
  Arrays can only be initialised \textit{at the time of declaration}.
  \begin{lstlisting}
    int arr[3] = {1, 2, 3};
  \end{lstlisting}
\end{remark}
As a function parameter, arrays can be taken as both a pointer or an array.
\begin{lstlisting}
  int sumArray(int *arr, int size) { ... }
  int sumArray(int arr[], int size) { ... }
\end{lstlisting}
Since the array name is a pointer, the function is able to modify the contents
of the passed array.

\paragraph{Multi-Dimensional Arrays} Given a array of type \verb|T|:
\[\verb|T a[D1][D2]|\ldots\verb|[Dn];|,\] 
where \verb|Dj| is the size of dimension $j$. $A_0$ is the address of the first
element \verb|a[0][0]|\ldots\verb|[0]|, and $S$ the size of each element of \verb|T|. We have
the address of the element \verb|a[i1][i2]|\ldots\verb|[in]| given by
\[A_0+S\times \left( \left( \left( \left( \verb|i1|\times \verb|D2|+\verb|i2| \right) \times \verb|D3|+\verb|i3| \right)\cdots \right)\times \verb|Dn|+\verb|in| \right) .\] 

\subsubsection{Strings}
A string is an \textit{array of characters terminated by a null character}
`\textbackslash0'. It can be defined by assigning each character to an array of
characters, or initialising the string directly:
\begin{multicols}{2}
  \begin{lstlisting}
    char str[6];
    str[0] = 'e';
    str[1] = 'g';
    str[2] = 'g';
    str[3] = '\0';
  \end{lstlisting}
  \begin{lstlisting}
    char str[] = "egg";
    char str[] = {'e', 'g', 'g', '\0'}
  \end{lstlisting}
  \vfill\null
\end{multicols}
Similarly, the string itself is a pointer referring to the address of the first
character of the string.

\begin{remark}
  We may include \verb|<string.h>| for string manipulation functions:
    \begin{itemize}
      \item \verb|int strlen(char *s)|
      \item \verb|int strcmp(char *s1, char *s2)|: s1 $-$ s2
      \item \verb|int strncmp(char *s1, char *s2, int n)|
      \item \verb|char *strcpy(char *s1, char *s2)|: s2 $\to$ s1
      \item \verb|char *strncpy(char *s1, char *s2, int n)|
    \end{itemize}
  It is typically preferred to use the functions with predefined maximum
  lengths $n$ to prevent exploits.
\end{remark}

\paragraph{String I/O}
In addition to the usual functions \verb|scanf("%s", str)| and
\verb|printf("%s\n", str)|,
strings can be read and printed using the\\
\centering
\verb|fgets(str, size, stdin)| and \verb|puts(str)| functions.\\
\raggedright
\begin{remark}
  \verb|fgets| also reads in the newline character, which may need to be
  replaced with the null character.
\end{remark}

\subsubsection{Structures}
\textit{Structure types} allow for grouping of heterogeneous members, and even
of other groups themselves.
\begin{lstlisting}
  typedef struct {
    type_1 x;
    type_2 y;
  } type_name;
\end{lstlisting}
\begin{remark}
  Structure types are definitions of types, not a declaration of a variable. No
  memory is allocated until a variable is declared.
\end{remark}
\begin{example}[Nested Structure Variable]
  Structure variables can then be initialised similarly to array
  initialisation, and members accessed using the dot \verb|.| operator:
  \begin{multicols}{2}
    \begin{lstlisting}
    typedef struct {
      int day, month, year;
    } date_t;

    typedef struct {
      int cardNum;
      date_t expiry;
    } card_t;
    \end{lstlisting}
    \begin{lstlisting}
      card_t card = {12345, {7, 10, 2025}};
      card.expiry.year = 2026;
      scanf("%d", &card.cardNum);

      // Assignments with structures are legal.
      card_t card_1 = {1, {1, 1, 1}};
      card_1 = card;
    \end{lstlisting}
    \vfill\null
  \end{multicols}
\end{example}
\begin{remark}
  Due to \textit{pass-by-value}, structures, when passed into a function, are
  \textbf{copied} into the formal parameters.
  To modify structure variables, the address of the variable is passed instead.
\end{remark}
We define the arrow operator \verb|->| as a ``shortcut'' syntax to access
members of a pointer to a structure variable:
\[\verb|(*struct_ptr).member|\equiv\verb|struct_ptr->member|\] 

\subsubsection{Unions and Enum}
Unions are a special type of structures in which the fields occupy the same
memory storage (that of the field requiring the greatest memory storage). That
is, only one of the fields can be defined in a union for defined behaviour.

Enum is a similar type that represents a group of integer constants, to
constrain options and choices.
